\section{Comparison of the three local routes}
\label{sec:route-comparison}
% Status: open. Compares route kernels, assignments, branch maps, scale placement, and compatibility tests.

The endpoint, interval, and \(G_2\) sections introduce three ways to place the DeVries determinant inside a theory with localized gauge and order-parameter data.  The section casts the three routes in a common reduction format.  Each route is tested by the same normalized two-channel kernel, electroweak assignment, branch interpretation, and global compatibility conditions.  A derivation in any route must supply a physical origin for the matrix entries and turn the string, extra-dimensional, or singular-geometric address into a calculation.

\subsection{Common reduction problem}
\label{sec:common-reduction-problem}

Let \(\mathcal S_r\) denote one of the three route descriptions,
\begin{equation}
  r\in
  \{\mathrm{end},\mathrm{int},G_2\},
  \qquad
  \mathcal S_r
  =
  \begin{cases}
  \mathcal S_{\rm end} & \hbox{open-string endpoint or D-brane sector},\\
  \mathcal S_{\rm int} & \hbox{interval Kaluza--Klein boundary sector},\\
  \mathcal S_{G_2} & \hbox{localized singularity sector}.
  \end{cases}
  \label{eq:route-source-list}
\end{equation}
The desired reduction chooses a two-dimensional light subspace
\begin{equation}
  \mathcal H_{J,r}=\operatorname{span}\{h_{J,r},a_{J,r}\},
  \label{eq:route-light-space}
\end{equation}
where \(h_{J,r}\) is the order-parameter channel and \(a_{J,r}\) is the adjoint/current channel.  In this basis the reduced spectral kernel has the form
\begin{equation}
  K_{J,r}(\lambda)=
  \begin{pmatrix}
  \lambda+\Sigma_{hh,J}^{r}(\lambda) &
  \Sigma_{ha,J}^{r}(\lambda)\\
  \Sigma_{ah,J}^{r}(\lambda) &
  \lambda+\Sigma_{aa,J}^{r}(\lambda)
  \end{pmatrix}.
  \label{eq:route-common-kernel}
\end{equation}
The DeVries target is the simultaneous limit and normalization
\begin{equation}
  \Sigma_{hh,J}^{r}=0,\qquad
  \Sigma_{aa,J}^{r}=J,\qquad
  \Sigma_{ha,J}^{r}\Sigma_{ah,J}^{r}=J.
  \label{eq:route-common-target}
\end{equation}
The available source equations select a candidate entry sequence.  The
interval/CHM route begins with the admissible neutral trace package of
Eqs.~\eqref{eq:chm-admissible-positive-norms}--\eqref{eq:chm-current-Z-factor}.
That package supplies positive projected norms, photon-orthogonal
normalization, source-ordered or commuting projectors, photon-reference
subtraction, a shared \(\Lambda_{\rm CHM}\), and the current-slope datum
\(Z_J^{\rm cur}\).  Its first theorem target is
\begin{equation}
  \widehat\Sigma_{aa,2}^{\rm CHM,can}=2 .
  \label{eq:route-chm-current-entry-first-test}
\end{equation}
Sec.~\ref{sec:kk-boundary} states the corresponding acceptance criteria in
terms of the CHM source equations and projector data.
The same source package then carries
\begin{equation}
  \Delta_W^{\rm CHM},\qquad
  \Delta_Y^{\rm CHM},\qquad
  \Sigma_{ha,J}\Sigma_{ah,J}=J,\qquad
  \Delta_{\rm O1/O10}.
  \label{eq:route-chm-followup-tests}
\end{equation}
The Hodge/SUSY-QM route supplies the companion square-root theorem target.  The
\(G_2\) route provides Target IIIg, a local ADE-pairing current test whose first
acceptance equation is
\(\widehat\Sigma_{aa,2}^{G_2,{\rm can}}=2\), and Target IIIh, a local
flow-overlap product target for
\(\widehat\Sigma_{ha,J}^{G_2}\widehat\Sigma_{ah,J}^{G_2}=J\).  Its
electroweak projectors, charge lattice, and pole map remain open obligations.  The
endpoint/Chan--Paton route
supplies a matrix arena whose DeVries use requires a worldvolume quadratic
kernel, projection, and normalization.
Equation~\eqref{eq:route-common-target} is stronger than a qualitative mechanism.  It assigns one geometric, endpoint, or boundary object to each matrix entry.  It also fixes the algebraic route to
\begin{equation}
  \det K_{J,r}(\lambda)=\lambda^2+J\lambda-J,
  \qquad
  \lambda\in\{x_+(J),x_-(J)\}.
  \label{eq:route-common-det}
\end{equation}

The route is fully specified when one source package supplies the reduced
space, source product, entry kernels, scalar functional, and pole map, and
derives the route-indexed arrows in
\begin{equation}
\begin{array}{ccc}
(u_r,\mathcal H_{J,r}) &
\longrightarrow &
(\langle\cdot,\cdot\rangle_r,\Lambda_r,
P_{h,J}^r,P_{a,J}^r,K_{\rm cur}^r,K_{\rm mix}^r)\\
\downarrow && \downarrow\\
(\mathcal F_{\rm sc}^r,\mathcal R_{\rm pole}^r)
& \longrightarrow &
(\widehat\Sigma_{aa,J}^r,
\widehat\Sigma_{ha,J}^r\widehat\Sigma_{ah,J}^r,
P_W^r,P_Z^r,P_\gamma^r,\Delta_{\rm match}^{(r)}) .
\end{array}
\label{eq:active-proof-spine}
\end{equation}
The current proof obligations include the canonical CHM neutral-current entry
\(\widehat\Sigma_{aa,2}^{\rm CHM,can}=2\), the \(G_2\) ADE current entry
\(\widehat\Sigma_{aa,2}^{G_2,{\rm can}}=2\), the same-basis Hodge/SUSY-QM or
\(G_2\) flow-overlap product
\(\widehat\Sigma_{ha,J}^r\widehat\Sigma_{ah,J}^r=J\), and the ordered
projector package \(P_W^r,P_Z^r,P_\gamma^r\) with its pole remainder.

The reduction has three candidate arenas.  The endpoint arena is supported by the open-string fact that endpoint labels and coincident branes generate matrix-valued gauge fields, with brane positions supplying scalar data \cite{TongString,Polchinski1994}.  The interval arena is supported by extra-dimensional models in which boundary conditions and brane-localized terms set vector spectra \cite{CsakiHubiszMeade2005}.  The \(G_2\) arena is supported by singular compactifications that localize gauge sectors and chiral matter, together with Higgs-bundle variables organizing local matter and interaction data \cite{WittenG2Anomaly2001,AcharyaWitten2001,AcharyaGukov2004,BraunEtAl2019}.  These sources support the listed ingredients.  Selection of the two-channel subspace, inner product, normalization, and entries in Eq.~\eqref{eq:route-common-target} remains the derivation target.

The interval arena supplies source-derived boundary data.  The CHM source
equations fix the boundary kinetic normalization, the corresponding modified
product, vector Robin data from boundary scalar vevs, scalar \(A_5/\pi_i\)
boundary equations, and photon/custodial bookkeeping.  Sec.~\ref{sec:kk-boundary}
states the convention data.  In compact route notation,
\begin{equation}
  \partial_y f_\lambda=-(\lambda/M_i)f_\lambda,\qquad
  (f,g)_{\rm CHM}=\int fg+\sum_i M_i^{-1}f(y_i)g(y_i),
  \qquad
  \partial_y A_\mu\mp v_i^2A_\mu=0.
  \label{eq:route-chm-source-data}
\end{equation}
Thus the source-backed statement is
\begin{equation}
  \hbox{CHM boundary data}
  \Longrightarrow
  K_{\rm int}(\lambda)\ \hbox{exists},
  \qquad
  K_{\rm int}(\lambda)
  \Longrightarrow
  \begin{pmatrix}
  \lambda&-\sqrt J\\
  -\sqrt J&\lambda+J
  \end{pmatrix}
  \quad\hbox{is open}.
  \label{eq:route-chm-open-target}
\end{equation}
Equation~\eqref{eq:route-chm-source-data} supplies boundary ingredients; the
projection to \(\mathcal H_J\) and the \(J\)-normalization are theorem data.
The interval route has explicit vector and scalar boundary equations for an
entry-level source target.  The entry tests have two linked forms:
Hodge/SUSY-QM factorization tests
\(\Sigma_{ha,J}\Sigma_{ah,J}=J\), and the canonical CHM boundary-current test
\(\widehat\Sigma_{aa,2}^{\rm CHM,can}=2\).  Sec.~\ref{sec:kk-schur-target}
defines the neutral trace package, photon subtraction, \(Z_J^{\rm cur}\),
canonical current kernel, and target equation.  The CHM equations supply the
boundary variation, product, vector Robin data, scalar \(A_5/\pi_i\) equations,
custodial bookkeeping, and photon-zero accounting.  The projection
\(P_{a,J}\), reference subtraction, source scale \(\Lambda_{\rm CHM}\),
endpoint-normal conventions, ordered W/Z boundary map, and pole map remain
theorem data.  The parent KK-fixing notes sharpen the same issue into a single-source
test: the dimensional middle line, negative-branch scalar reading, and
electromagnetic endpoint can be combined when one source variable
supplies
\begin{equation}
  u\mapsto
  \big(t_{\rm dim}(u),t_{\rm EW}(u),
  K_J(u,\lambda),\mathcal F_{\rm sc}(u;J)\big).
  \label{eq:route-single-source-chm-test}
\end{equation}
The scalar and electromagnetic readings are downstream of the current entry,
ordered W/Z assignment, source-to-pole calculation, and photon-zero
normalization.  Target~\ref{app:target-alpha-endpoint} records the O17 endpoint
coupling package after those data are supplied.

The common reduction is filtered by three admissibility tests.  First, trace
and coefficient statements stay inside the matrix where they are derived until
a common source operator and field normalization identify the trace spaces.
Second, Wigner--Eckart or Clebsch routes must specify the parent
representation, operator tensor type, projection, and normalization.  The
tested doublet-parent route is removed by the ordinary selection rule.  The
triplet axial route gives the wrong parent-Casimir datum for
Eq.~\eqref{eq:route-common-target}.  Third, an SO(32)-flavor completion must
show how its representation sector couples to the electroweak kernel, with
extra-state projection and anomaly constraints stated before any flavor claim
is used in the mass-ratio argument.

For the superconnection/CHM route, these filters have concrete meanings.  The
odd Higgs trace space, the even transverse-current trace space, and the
interval boundary scalar product stay separated until one source operator
defines their common inner product.  The Wigner--Eckart route is a failed
subroute, so the superconnection route must supply its own projection.  The
SO(32) material remains flavor-boundary bookkeeping until a
derived coupling inserts that representation sector into the electroweak kernel.

\subsection{Entry-by-entry comparison}
\label{sec:entry-comparison}

The first entry is the diagonal current contribution,
\begin{equation}
  a_{J,r}
  \quad\leadsto\quad
  \Sigma_{aa,J}^{r}=J.
  \label{eq:route-current-question}
\end{equation}
The candidate source objects are a Chan--Paton or brane gauge amplitude, a
boundary value of a bulk gauge mode, or a fluctuation along a singular gauge
locus.

The second entry is the off-diagonal product,
\begin{equation}
  h_{J,r}
  \leftrightarrow
  a_{J,r},
  \qquad
  \Sigma_{ha,J}^{r}\Sigma_{ah,J}^{r}=J.
  \label{eq:route-mixing-question}
\end{equation}
The candidate source objects are a stretched open string or endpoint overlap,
a boundary overlap between an order parameter and a bulk gauge profile, or a
localized matter/deformation coupling.

The third entry is the order-parameter reference contribution,
\begin{equation}
  h_{J,r}
  \quad\leadsto\quad
  \Sigma_{hh,J}^{r}=0.
  \label{eq:route-h-reference}
\end{equation}
The candidate source objects are a boundary reference channel, a normalized
boundary order-parameter self-energy, or a critical-point normalization of a
localized mode.  Appendix~\ref{app:target-boundary-determinant} records the
entry-by-entry proof obligations.

The following table records the three entries in a common dictionary.
\begin{center}
\small
\begin{tabular}{p{0.18\linewidth}p{0.24\linewidth}p{0.24\linewidth}p{0.24\linewidth}}
\toprule
Kernel entry & Endpoint/D-brane route & Interval KK route & \(G_2\) route\\
\midrule
\(\Sigma_{aa}=J\) & Chan--Paton current or brane gauge amplitude & CHM boundary current-kernel target with source conventions; \(J\)-normalization open & Target IIIg: \(A_1\subset A_2\) ADE current pairing with \(\widehat\Sigma_{aa,2}^{G_2,{\rm can}}=2\)\\
\(\Sigma_{ha}\Sigma_{ah}=J\) & stretched string or endpoint overlap & boundary scalar vev, \(A_5/\pi_i\) equation, and modified product; open product normalization & Target IIIh: localized matter at \(df_Q=0\), oriented flow kernel, bilinear maps, and point-support metric\\
\(\Sigma_{hh}=0\) & order-parameter reference channel & boundary scalar reference term; open subtraction rule & critical-point normalization with local-to-pole matching\\
\bottomrule
\end{tabular}
\end{center}

\subsection{Electroweak assignment as a route-independent target}
\label{sec:route-ew-assignment}

The same electroweak assignment target applies to every route.  The low-energy field-theory data contain an \(SU(2)_L\) Higgs doublet and adjoint gauge currents.  Their available quadratic invariants give
\begin{equation}
  J_H=\frac34,\qquad J_{\rm adj}=2.
  \label{eq:route-j-invariants}
\end{equation}
The target map is
\begin{equation}
  \mathcal A_r:
  (\mathcal S_r,\mathcal H_{J,r})
  \longrightarrow
  (J_W,J_Z)=(J_H,J_{\rm adj}).
  \label{eq:route-assignment-map}
\end{equation}
The ordered assignment is a theorem target.  The observed W and Z vector bosons
are gauge-sector particles after symmetry breaking.  The quotient assigns the
charged comparison to the order-parameter invariant and the neutral comparison
to the adjoint/current invariant.  The required derivation must produce this
ordered pair from the gauge-Higgs coupling structure, endpoint boundary labels
attached to the Higgsing field, interval boundary data that distinguish charged
and neutral sectors, or a localized singularity model where doublet and adjoint
data meet.
The Standard Model mass matrix supplies the projective ratio and the
representation data; the ordered sampling maps remain additional theorem data.

If the ordered assignment and pole-placement theorem hold, the pole-ratio
reading has one route-labeled form:
\begin{equation}
  \frac{M_{W,\rm pole}^2}{M_{Z,\rm pole}^2}
  =
  \frac{x_+(J_H)}{x_+(J_{\rm adj})}
  +\Delta^{(r)}_{\rm match}.
  \label{eq:route-matching-remainder}
\end{equation}
Here \(r\) labels the endpoint, interval, or \(G_2\) route.  The exact pole
reading requires a derived \(\Delta^{(r)}_{\rm match}=0\) in the chosen pole
scheme, or a computed remainder whose size and sign are fixed by the
source-to-EFT-to-pole chain.  A compactification-scale or unification-scale
derivation makes the quotient a boundary datum; the pole claim then requires
the matching chain.

\subsection{Branch interpretation and scalar data}
\label{sec:route-branch-interpretation}

Every route diagonalizes the same two-channel kernel.  Therefore each route has to account for both eigenvectors.  Let
\begin{equation}
  u_{\pm,J,r}=
  \alpha_{\pm,J,r}h_{J,r}
  +
  \beta_{\pm,J,r}a_{J,r}
  \label{eq:route-eigenvectors}
\end{equation}
be normalized eigenvectors of the route kernel in the DeVries limit.  The positive branch supplies the pole-ratio clue.  The negative branch supplies the scalar-side question,
\begin{equation}
  u_{-,J,r}
  \quad\leadsto\quad
  \mathcal F_{{\rm sc},r}.
  \label{eq:route-scalar-map}
\end{equation}
The functional \(\mathcal F_{{\rm sc},r}\) depends on the route.  Endpoint, interval, and \(G_2\) sources provide candidate scalar addresses: brane or boundary moduli, localized scalar modes, and singularity-deformation data.  These addresses acquire physical content after a gauge-invariant functional, normalization, and electroweak scheme are specified.
Appendix~\ref{app:target-negative-branch} states the same-source condition:
the first five entries of the scalar package must agree with the vector pole
package, including \(\mathcal H_J\), the inner product, the projection,
\(\Lambda_J\), and \(\mathcal R_{\rm pole}\).

The scalar map has to be tied to the electroweak ray.  The gauge-Higgs system has a radial vacuum scale and a projective vector-mass ratio.  The two-root structure is physically restrictive because a determinant used for the pole clue also supplies the partner branch.  A source route must identify how the positive branch fixes the projective vector datum and how the negative branch enters the radial or scalar sector.

\subsection{Scale placement}
\label{sec:route-scale-placement}

The route comparison separates three scale placements:
\begin{equation}
  \hbox{pole scale},
  \qquad
  \hbox{compactification or boundary scale},
  \qquad
  \hbox{unification scale}.
  \label{eq:route-scale-placements}
\end{equation}
The construction attaches the electroweak clue to the pole scale.  The
attachment gives the clue its direct electroweak content.  The GUT lineage
supplies the comparison class: a clean weak-angle value can be a high-scale
structural datum, then acquire low-energy meaning through running and
thresholds.  In the present construction the clean value is attached to the pole
ratio, so the scale-placement theorem is logically prior to any compactification
geometry.

The common matching theorem has the form
\begin{equation}
  \mathcal S_r
  \longrightarrow
  K_{J,r}^{\rm bare}(\lambda)
  \longrightarrow
  \Gamma_{\rm eff}^{(4)}[H,W,B]
  \longrightarrow
  \Delta^{-1}_{V,T}(s)
  \big|_{s=s_{W,Z}} .
  \label{eq:route-pole-matching-chain}
\end{equation}
Equation~\eqref{eq:route-pole-matching-chain} is the pole-placement theorem
target.  The first arrow is the route reduction.  The second arrow matches the source
kernel into a four-dimensional gauge-Higgs effective action.  The third arrow
places the result in the dressed transverse pole convention of Sec.~\ref{sec:pole}.
Endpoint, interval, and \(G_2\) routes differ in the first arrow; they share
the last two arrows when the pole reading is retained.
Appendix~\ref{app:target-pole-placement} gives the equivalent transverse
self-energy form of the same theorem.

For endpoint and interval routes, the central obligation is the relation between
a boundary determinant and a four-dimensional pole.  Boundary conditions define
a spectrum before loop corrections and before conversion to pole observables.
The derivation specifies how the determinant survives matching to the pole
definition, or how the pole reading emerges from a dressed boundary
self-energy.

For the \(G_2\) route, a local singularity model lives at a compactification or
high scale.  The local kernel can organize representations and couplings.  A
pole-spectrum reading requires a low-energy matching chain.  The resulting
target is a two-stage theorem:
\begin{equation}
  \mathcal S_{G_2}
  \longrightarrow
  K_{J,G_2}
  \longrightarrow
  \Delta^{-1}_{V,T}(s)
  \quad\hbox{with}\quad
  s=s_{W,Z}.
  \label{eq:g2-two-stage-pole}
\end{equation}
The first arrow is local geometry.  The second arrow is four-dimensional matching.  Both arrows remain open obligations.

\subsection{Route outcome}
\label{sec:route-outcome}

A route succeeds when it supplies the common kernel target, the ordered
electroweak assignment, the pole-matching chain, the scalar-branch map, and
the global compatibility checks from one stated source theory.  These criteria
turn endpoint, interval, and \(G_2\) language into comparable derivation
problems.  Extra light channels, a different trace, a different mixing product,
a high-scale placement, or an incompatible global form would classify the
DeVries relation inside a larger family of spectral constructions.
